\label{sec:prong3}

\subsection{The Legal Standard}

Prong 3 requires that the majority ``vote sufficiently as a bloc to enable it
\ldots\ usually to defeat the minority's preferred candidate.''  \textit{Gingles},
478 U.S. at 51.  The key quantities are: (a) the majority vote share for the
candidate opposed by the minority-preferred candidate; and (b) the frequency
with which majority bloc voting defeats the minority-preferred candidate
across a sample of elections.

\subsection{The \textit{Callais} Disentanglement Requirement}

\textit{Louisiana v.\ Callais}~\citep{callais2026} added a methodological
requirement that has become the central battleground in post-2026 Section 2
litigation: expert witnesses must disentangle racial bloc voting from partisan
bloc voting in Prong 3 analysis.  The disentanglement requirement responds
to a specific litigation strategy: states defend against Prong 3 findings by
arguing that observed majority bloc voting is partisan (white Republicans
voting against Black Democratic candidates), not racial (white voters voting
against Black candidates), and is therefore constitutionally permissible.

The \textit{Callais} Court held that this is a substantive question requiring
expert analysis, not a bare assertion.  An expert witness who simply reports
that majority and minority voters vote differently without controlling for
party has not satisfied Prong 3 under the post-\textit{Callais} standard.

\subsection{The WLS+HC3+Holm Regression}

The post-\textit{Callais} Prong 3 analysis uses a two-equation system.
The first equation estimates majority racial cohesion with a partisan control:

\begin{equation}
  Y_t = \alpha + \beta_{\text{race}} \cdot X^{\text{maj}}_t
                + \beta_{\text{party}} \cdot P_t
                + \varepsilon_t
  \label{eq:prong3-main}
\end{equation}

where $P_t$ is the precinct-level Republican presidential vote share, which
serves as the partisan control variable.  The estimator $\hat{\beta}_{\text{race}}$
isolates the racial component of majority bloc voting---the portion not
attributable to partisan alignment.

The regression is estimated by WLS with weights proportional to precinct
turnout.  Standard errors are computed using the HC3 heteroskedasticity-consistent
covariance matrix estimator~\citep{mackinnon1985some}, which is robust to
the heteroskedasticity present in precinct-level election data.

\subsection{Multiple Election Adjustment with Holm Correction}

Prong 3 is typically analyzed across multiple elections to demonstrate that
majority bloc voting is a persistent pattern.  Testing across multiple
elections creates a multiple-testing problem: with enough elections, a
researcher will find a statistically significant bloc-voting result by chance.
The \textit{Callais} decision implies that courts should be skeptical of
Prong 3 findings based on a single election or on uncorrected multiple-testing
results.

We apply the Holm sequential rejection procedure~\citep{holm1979simple} to
correct for multiple testing across elections.  The Holm procedure is more
powerful than the Bonferroni correction while controlling the family-wise
error rate.  The corrected $p$-values are reported alongside the uncorrected
values to allow the court to assess the sensitivity of the finding to the
multiple-testing correction.

\subsection{LOO Robustness Checks}

Following the implementation in \texttt{redist-analysis::bloc\_voting}, we
run leave-one-out (LOO) variants: for each election in the sample, we refit
the regression excluding that election and verify that the Prong 3 conclusion
is stable.  An unstable conclusion---one that reverses when a single election
is excluded---would indicate that the result is driven by a single unusual
election rather than a structural pattern.

\subsection{Running the Prong 3 Analysis}

The full post-\textit{Callais} Prong 3 analysis is produced by:
\begin{verbatim}
redist analyze --state AL --year 2020 --version remedy \
  --types bloc-voting --partisan-control
\end{verbatim}
The \texttt{--partisan-control} flag activates the partisan control variable
in Equation~\ref{eq:prong3-main}, the HC3 standard errors, the Holm correction
across elections, and the LOO robustness checks.  The output includes: (a) the
racial bloc-voting coefficient $\hat{\beta}_{\text{race}}$ for each election;
(b) the Holm-corrected $p$-values; (c) the LOO stability table; and (d) a
narrative summary of the Prong 3 findings suitable for inclusion in an expert
report.
